

\documentclass[12pt,letterpaper]{article}

\usepackage[utf8]{inputenc}
\usepackage[T1]{fontenc}
\usepackage[spanish,es-noquoting,es-tabla]{babel}
\usepackage{carlito}                 % Calibri
\usepackage[margin=2.5cm]{geometry}
\usepackage{setspace}
\onehalfspacing                      % interlineado 1,5
\usepackage{tabularx}
\usepackage{booktabs}
\usepackage{array}
\usepackage{float}
\usepackage{enumitem}
\usepackage[most]{tcolorbox}
\usepackage[hidelinks]{hyperref}

\setlength{\parskip}{0.4em}
\setlength{\parindent}{0pt}
\setlist{nosep, leftmargin=1.4em}
\renewcommand{\arraystretch}{1.15}
\newcolumntype{L}{>{\raggedright\arraybackslash}X}
\newcolumntype{C}[1]{>{\centering\arraybackslash}p{#1}}

% ---------------------------------------------------------------------
% >>> Integrantes
\newcommand{\IntA}{Carlos Alfonso Gonzalez Perez}
\newcommand{\IntB}{Carlos Sebastian Madrigal Rodriguez}
\newcommand{\IntC}{3}
% ---------------------------------------------------------------------

\newtcolorbox{redactar}[1]{
  colback=gray!6, colframe=gray!45, boxrule=0.5pt, arc=2pt,
  left=6pt, right=6pt, top=4pt, bottom=4pt,
  fonttitle=\bfseries\small, coltitle=black, colbacktitle=gray!18,
  title={#1}, before skip=6pt, after skip=8pt
}

\begin{document}

\begin{center}
  {\Large\bfseries Planeación de un sprint en el marco de trabajo Scrum}\\[0.3em]
  {\large Portal de solicitudes de vacaciones con flujo de aprobación}\\[0.6em]
  \IntA \quad·\quad \IntB \quad·\quad \IntC\\[0.2em]
  {\small\today}
\end{center}

% =====================================================================
\section{Descripción del proyecto}

\textbf{Empresa.} Nexora Soluciones es una consultora de software que se especializa en aplicaciones para gestionar procesos internos. \textbf{Cliente.} Su cliente, Grupo Industrial Valtek, tiene 800 empleados y actualmente maneja las vacaciones con formularios en papel, lo que causa problemas en la gestión.

\textbf{Producto.} Un portal web permite a cada empleado consultar su saldo de días, registrar solicitudes de vacaciones y seguir su estado. Los jefes aprueban o rechazan y Recursos Humanos obtiene reportes.

\textbf{Objetivo.} Reducir el tiempo de aprobación de una solicitud de más de siete días a menos de 48 horas y eliminar los errores en el cálculo de saldos, dejando trazabilidad completa de cada decisión.

% =====================================================================
\section{Asignación de roles Scrum}

Al ser tres integrantes, se combinan roles respetando que el product owner no
desarrolle, para que pueda representar al cliente sin conflicto de intereses.

\begin{table}[H]
\centering
\caption{Roles del equipo Scrum.}
\begin{tabularx}{\textwidth}{lL}
\toprule
\textbf{Integrante} & \textbf{Rol y responsabilidad principal} \\
\midrule
\IntA & \textbf{Product owner.} Representa al cliente, prioriza el
        \emph{product backlog} y valida los criterios de aceptación. \\
\IntB & \textbf{Scrum master y desarrollador.} Facilita los eventos,
        elimina impedimentos y desarrolla la capa de servidor. \\
\IntC & \textbf{Desarrollador.} Interfaz web, pruebas y despliegue. \\
\bottomrule
\end{tabularx}
\end{table}

% =====================================================================
\section{Product backlog}

Esfuerzo estimado en puntos de historia (escala de Fibonacci: 1, 2, 3, 5, 8,
13), donde un valor mayor indica más complejidad e incertidumbre.

\begin{table}[H]
\centering
\caption{Product backlog priorizado.}
\small
\begin{tabularx}{\textwidth}{lLcc}
\toprule
\textbf{ID} & \textbf{Historia de usuario} & \textbf{Prioridad} & \textbf{Puntos} \\
\midrule
HU-01 & Como \textbf{empleado}, quiero iniciar sesión con mi número de
        empleado para acceder a mis datos de forma segura. & Alta & 3 \\
HU-02 & Como \textbf{empleado}, quiero consultar mi saldo de días disponibles
        para planear mis vacaciones. & Alta & 3 \\
HU-03 & Como \textbf{empleado}, quiero registrar una solicitud indicando
        fechas para pedir mis vacaciones sin usar papel. & Alta & 5 \\
HU-04 & Como \textbf{jefe inmediato}, quiero aprobar o rechazar solicitudes
        con un comentario para responder en tiempo. & Alta & 5 \\
HU-05 & Como \textbf{empleado}, quiero recibir un correo cuando cambie el
        estado de mi solicitud para no tener que revisar el portal. & Media & 3 \\
HU-06 & Como \textbf{analista de RR.\,HH.}, quiero un reporte de vacaciones por
        planta y departamento para planear la cobertura. & Media & 8 \\
HU-07 & Como \textbf{analista de RR.\,HH.}, quiero consultar el historial de
        cada solicitud para auditar quién la aprobó y cuándo. & Baja & 5 \\
\bottomrule
\end{tabularx}
\end{table}

% =====================================================================
\section{Sprint backlog}

\textbf{Objetivo del sprint (una semana).} Un empleado puede revisar su saldo y enviar una solicitud, mientras que su jefe puede aprobarla o rechazarla. Se escogen las HU-02, HU-03 y HU-04, que componen un flujo completo para prueba, sumando 13 puntos, factibles para dos desarrolladores en cinco días.

\begin{table}[H]
\centering
\caption{Desglose de tareas del sprint.}
\small
\begin{tabularx}{\textwidth}{lLlc}
\toprule
\textbf{HU} & \textbf{Tarea} & \textbf{Responsable} & \textbf{Horas} \\
\midrule
HU-02 & Diseñar la tabla de saldos y la regla de cálculo por antigüedad & \IntB & 3 \\
      & Programar el servicio que devuelve el saldo del empleado          & \IntB & 4 \\
      & Maquetar la vista de saldo en el portal                            & \IntC & 3 \\
      & Validar el cálculo con casos reales del cliente                    & \IntA & 2 \\
\midrule
HU-03 & Programar el registro de solicitudes y la validación de fechas    & \IntB & 6 \\
      & Construir el formulario con calendario y mensajes de error         & \IntC & 5 \\
      & Escribir pruebas de fechas traslapadas y saldo insuficiente        & \IntC & 3 \\
      & Revisar los criterios de aceptación de la historia                 & \IntA & 1 \\
\midrule
HU-04 & Programar la bandeja de pendientes y la acción de aprobar/rechazar & \IntB & 5 \\
      & Construir la vista del jefe con filtro por estado                  & \IntC & 4 \\
      & Desplegar en el entorno de pruebas y preparar la demo              & \IntC & 2 \\
      & Validar el flujo completo antes de la \emph{sprint review}         & \IntA & 2 \\
\midrule
      & \textbf{Total}                                                     &       & \textbf{40} \\
\bottomrule
\end{tabularx}
\end{table}

% =====================================================================
\section{Reflexión individual}


\subsection*{\IntA\ — Product owner}

\begin{redactar}
\small
Durante el sprint, como encargado del producto, mi función principal sería garantizar que el grupo comprenda las demandas del cliente y que las historias elegidas ofrezcan valor al producto. Me ocuparía de evaluar y organizar el Product Backlog, aclarar inquietudes sobre los requisitos y verificar que las funcionalidades creadas satisfagan los criterios de aceptación establecidos. Uno de los principales obstáculos sería mantener una comunicación continua con el equipo y prevenir modificaciones innecesarias que puedan impactar en el objetivo del sprint. También prevendría como desafío definir adecuadamente las normas para el cálculo de saldos y la validación de las solicitudes, puesto que cualquier falla podría repercutir directamente en el usuario. Aportaría al esfuerzo del equipo brindando información pertinente, constatando los avances y participando en las revisiones. Considero que mi involucramiento resultaría esencial para mantener el énfasis en las necesidades del cliente y asegurar que el resultado final sea útil, claro y práctico.
\end{redactar}

\subsection*{\IntB\ — Scrum master y desarrollador}

\begin{redactar}
\small
Durante el sprint, como Scrum Master y desarrollador, facilitaría las reuniones
diarias para que el equipo comparta sus avances, identifique impedimentos y
mantenga el enfoque en el objetivo acordado. También apoyaría la preparación de
la revisión y la retrospectiva, procurando que la comunicación entre el Product
Owner y los desarrolladores sea constante y que cualquier ajuste se valore de
acuerdo con las prioridades del producto. En mi función de desarrollo,
implementaría los servicios de servidor necesarios para consultar los saldos,
registrar solicitudes y permitir que los jefes las aprueben o rechacen.

Como retos principales anticiparia el mantener a la gente interesada y motivada 
durante el sprint, asegurando que todos comprendan sus responsabilidades, se sientan 
motivados y alineados con los objetivos del equipo. Aparte del enfoque de equipo,
otro gran reto seria mantener los requerimientos claros y bien definidos, evitando 
malentendidos y cambios inesperados que puedan afectar el progreso del equipo.
Para dichos retos, aclararía los criterios de aceptación desde el inicio, dividiría 
el trabajo en tareas verificables y colaboraría con el equipo en las pruebas del 
flujo completo.
\end{redactar}

\subsection*{\IntC\ — Desarrollador}

\begin{redactar}
\small
¿Qué tareas realizarías en tu rol durante el sprint? · ¿Qué retos anticipas? ·
¿Cómo contribuirías al trabajo del equipo?
\end{redactar}

% =====================================================================
\section{Conclusión grupal}

\begin{redactar}
\small
¿Qué aprendimos como equipo sobre Scrum y el trabajo colaborativo? Resumir
también las decisiones tomadas (por ejemplo, por qué se eligieron esas tres
historias) y su impacto en el proceso.
\end{redactar}

\end{document}
